% ==================================================================
%  徐森林、金亚东、薛春华《数学分析》（第三册）
%  复习题 14
%  源：原书印刷页 172-177（PDF 物理页 178-183）
%  底本：PaddleOCR-VL-1.6，见 chapters/_ch14_14_3_raw.md 末尾
%  约定：只用 exblock，不设 section（故不入目录）
%  全 21 题题面照录；订正说明见 tools/_make_ch14_ex.py 头部
% ==================================================================

\begin{exblock}{复习题 14}
\begin{enumerate}
  \item 设 $\sum_{n=0}^{\infty}a_{n}$ 为一个发散的正项级数, 如果 $\lim\limits_{n\to+\infty}\frac{a_{n}}{a_{1}+a_{2}+\cdots+a_{n}}=0$, 证明: 幂级数 $\sum_{n=0}^{\infty}a_{n}x^{n}$ 的收敛半径 $R=1$.
  \item 设 $f(x)$ 及其所有导数在区间 $[0,r]$ 上都是非负的.

  (1) 应用 $f(x)$ 的 Taylor 展开式的积分余项
  \[
    R_{n}(x)=\frac{1}{n!}\int_{0}^{x}(x-t)^{n}f^{(n+1)}(t)\mathrm{d}t
\]
  证明:
  \[
    0\leqslant R_{n}(x)\leqslant\left(\frac{x}{r}\right)^{n+1}f(r),\quad x\in[0,r];
\]
  (2) 应用 (1) 证明: $f(x)$ 能在 $[0,r)$ 上展开为 Taylor 级数
  \[
    f(x)=\sum_{n=0}^{\infty}\frac{f^{(n)}(0)}{n!}x^{n};
\]
  (3) 应用 (2) 及例 14.2.8 以外的方法证明: $f(x)$ 能在 $[0,r)$ 上展开为 Taylor 级数.
  \item 设
  \[
    a_{n}=\begin{cases}
      \mathrm{C}_{3}^{k_{m}}, & n=3^{m}+k\ (m=1,2,\cdots;\ k=0,1,2,\cdots,3^{m}),\\
      0, & \text{其他},
\end{cases}
\]
  证明:

  (1) $\displaystyle\sum_{m=1}^{\infty}\left(\mathrm{C}_{3}^{0}x^{3^{m}}+\mathrm{C}_{3}^{1}x^{3^{m}+1}+\cdots+\mathrm{C}_{3}^{3^{m}}x^{3^{m}+3^{m}}\right)=\sum_{m=1}^{\infty}[x(1+x)]^{3^{m}}$;

  (2) $\sum_{n=1}^{\infty}a_{n}x^{n}$ 的收敛半径为 $\frac{\sqrt{5}-1}{2}$;

  (3) $\lim\limits_{n\to+\infty}\sqrt[n]{a_{n}}=\frac{\sqrt{5}-1}{2}$.
  \item 设 $S_{n}=\sum_{k=0}^{n}a_{k}$, $\sigma_{n}=\frac{S_{0}+S_{1}+\cdots+S_{n-1}}{n}$. 证明:

  (1) 如果 $\{\sigma_{n}\}$ 收敛, 则 $a_{n}=o(n)$, $n\to+\infty$;

  (2) 如果 $\{\sigma_{n}\}$ 收敛, 则 $f(x)=\sum_{n=0}^{\infty}a_{n}x^{n}$ 在 $(-1,1)$ 内绝对收敛, 且
  \[
    f(x)=(1-x)^{2}\sum_{n=0}^{\infty}(n+1)\sigma_{n+1}x^{n};
\]
  (3) 如果 $\lim\limits_{n\to+\infty}\sigma_{n}=S$, 则 $\lim\limits_{x\to1^{-}}f(x)=S$.
  \item 设 $f(x)$ 在 $[0,1]$ 上连续, 令
  \[
    f_{0}(x)=f(x),\quad f_{n+1}(x)=\int_{x}^{1}f_{n}(y)\mathrm{d}y,\quad x\in[0,1],\ n=0,1,2,\cdots.
\]
  证明:

  (1) $\sum_{n=1}^{\infty}f_{n}(x)$ 在 $[0,1]$ 上绝对一致收敛;

  (2) 记 $\varphi(x)=\sum_{n=1}^{\infty}f_{n}(x)$, 则
  \[
    \varphi'(x)+\varphi(x)+f(x)=0.
\]
  从而, $\sum_{n=1}^{\infty}f_{n}(x)$ 在 $[0,1]$ 上一致收敛于 $\varphi(x)=\int_{x}^{1}\mathrm{e}^{y-x}f(y)\mathrm{d}y$;

  (3) 将 $\varphi(x)=u(x)\mathrm{e}^{-x}$ 代入 $\varphi'(x)+\varphi(x)+f(x)=0$, 则得
  \[
    u(x)=\int_{x}^{1}f(y)\mathrm{e}^{y}\mathrm{d}y,\quad\varphi(x)=\int_{x}^{1}f(y)\mathrm{e}^{y-x}\mathrm{d}y.
\]
  \item 证明:

  (1) 幂级数 $S(x)=\displaystyle\sum_{k=1}^{\infty}\frac{(-1)^{k-1}}{3k-1}x^{3k-1}$ 的收敛点集为 $(-1,1]$;

  (2)
  \[
    \begin{aligned}
      S(x)&=\int_{0}^{x}S'(t)\mathrm{d}t=\int_{0}^{x}\frac{t}{1+t^{3}}\mathrm{d}t\\
      &=-\frac{1}{3}\ln(1+x)+\frac{1}{6}\ln(1-x+x^{2})+\frac{1}{\sqrt{3}}\arctan\left[\frac{2}{\sqrt{3}}\left(x-\frac{1}{2}\right)\right]+\frac{1}{\sqrt{3}}\arctan\frac{1}{\sqrt{3}}
    \end{aligned}
\]
  (3) $\displaystyle\frac{1}{2}-\frac{1}{5}+\frac{1}{8}-\frac{1}{11}+\cdots=\sum_{k=1}^{\infty}\frac{(-1)^{k-1}}{3k-1}=\frac{\sqrt{3}}{9}\pi-\frac{1}{3}\ln2$.
  \item 证明: $\displaystyle1-\frac{1}{4}+\frac{1}{7}-\frac{1}{10}+\cdots=\sum_{k=0}^{\infty}\frac{(-1)^{k}}{3k+1}=\frac{\sqrt{3}}{9}\pi+\frac{1}{3}\ln2$.
  \item 证明:

  (1) 幂级数
  \[
    S(x)=1+\sum_{n=1}^{\infty}\frac{(2n-1)!!}{(2n)!!}x^{n}
\]
  的收敛半径 $R=1$, 收敛点集 $[-1,1)$;

  (2) $2S'(x)-2xS'(x)=S(x)$; $S(x)=\dfrac{1}{\sqrt{1-x}}$, $x\in[-1,1)$;

  (3) $\displaystyle1+\sum_{n=1}^{\infty}(-1)^{n}\frac{(2n-1)!!}{(2n)!!}=\frac{\sqrt{2}}{2}$.
  \item (1) 用逐项积分求幂级数 $\displaystyle\sum_{n=1}^{\infty}(-1)^{n}\frac{n(n+1)}{2^{n}}x^{n-1}$ 的和;

  (2) 求极限 $\displaystyle\lim_{x\to0}\sum_{n=1}^{\infty}\frac{(-1)^{n}}{x^{2}+\frac{2^{n}}{n(n+1)}}$.
  \item 利用 $\sum_{n=1}^{\infty}\frac{1}{n^{2}}=\frac{\pi^{2}}{6}$, 证明:
  \[
    \int_{0}^{1}\frac{\ln x}{1-x^{2}}\mathrm{d}x=-\frac{\pi^{2}}{8}.
\]
  \item 设 $f(x)=\sum_{n=0}^{\infty}a_{n}x^{n}$ ($a_{n}>0$, $n=0,1,2,\cdots$) 的收敛半径为 $+\infty$, 且 $\sum_{n=0}^{\infty}a_{n}n!$ 收敛, 则 $\displaystyle\int_{0}^{+\infty}\mathrm{e}^{-x}f(x)\mathrm{d}x$ 也收敛, 且 $\displaystyle\int_{0}^{+\infty}\mathrm{e}^{-x}f(x)\mathrm{d}x=\sum_{n=0}^{\infty}a_{n}n!$.
  \item 利用已知公式
  \[
    \int_{-1}^{1}\frac{\mathrm{d}x}{(\alpha-x)\sqrt{1-x^{2}}}=\frac{\pi}{\sqrt{\alpha^{2}-1}},\quad\alpha>1
\]
  证明: $\displaystyle\int_{-1}^{1}\frac{x^{2n}}{\sqrt{1-x^{2}}}\mathrm{d}x=\frac{(2n-1)!!}{(2n)!!}\pi$.
  \item 设幂级数 $f(x)=a_{0}+a_{1}x+\cdots+a_{n}x^{n}+\cdots$ 在 $x=1$ 处收敛, 证明: 极限
  \[
    \lim_{v\to+\infty}\frac{v^{v}}{2^{v}v!}\left[f^{(n)}(0)-f^{(n)}\left(\frac{2^{v}v!}{v^{v}}\right)\right]=-a_{n+1}(n+1)!
\]
  \item 设 $f(x)$ 在 $(-\infty,+\infty)$ 上无穷次可导且满足:

  (1) $\exists M$, 使得 $|f^{(k)}(x)|\leqslant M$, $\forall x\in(-\infty,+\infty)$, $k=0,1,2,\cdots$;

  (2) $f\left(\frac{1}{2^{n}}\right)=0$, $n=1,2,\cdots$.

  证明: 在 $(-\infty,+\infty)$ 上, $f(x)=0$.
  \item 设 $f(x)=\displaystyle\sum_{n=1}^{\infty}\frac{x^{n}}{n^{2}}$, $x\in[0,1]$, 应用逐项求导与 $\sum_{n=1}^{\infty}\frac{1}{n^{2}}=\frac{\pi^{2}}{6}$ 证明:
  \[
    f(x)+f(1-x)+\ln x\cdot\ln(1-x)=\frac{\pi^{2}}{6},\quad x\in(0,1).
\]
  \item 证明: (Tauber 定理) 设在 $-1<x<1$ 上有
  \[
    f(x)=\sum_{n=0}^{\infty}a_{n}x^{n},\quad\lim_{n\to+\infty}na_{n}=0.
\]
  如果 $\lim\limits_{x\to1^{-}}f(x)=S$, 则 $\sum_{n=0}^{\infty}a_{n}$ 收敛且其和为 $S$.
  \item 设 $P_{n}(x)=1+x+\frac{x^{2}}{2!}+\cdots+\frac{x^{n}}{n!}$, $x_{m}$ 为 $P_{2m+1}(x)=0$ 的实根. 证明: $x_{m}<0$, 且 $\lim\limits_{m\to+\infty}x_{m}=-\infty$.
  \item 设 $\{f_{n}(x)\}$ 在任何有界区间上一致收敛于 $f(x)$, 且 $f_{n}(x)$ 在任何有界区间上可积, 又存在函数 $F(x)$, 使 $|f_{n}(x)|\leqslant F(x)$, $n=1,2,\cdots$, 且 $\displaystyle\int_{-\infty}^{+\infty}F(x)\mathrm{d}x$ 收敛. 证明:
  \[
    \lim_{n\to+\infty}\int_{-\infty}^{+\infty}f_{n}(x)\mathrm{d}x=\int_{-\infty}^{+\infty}f(x)\mathrm{d}x.
\]
  \item 设 $f_{n}(x)$ 在 $(-\infty,+\infty)$ 上一致连续, 且 $\{f_{n}(x)\}$ 在 $(-\infty,+\infty)$ 上一致收敛于 $f(x)$. 证明: $f(x)$ 在 $(-\infty,+\infty)$ 上也一致连续.
  \item 证明: 每个幂级数 $\sum_{n=0}^{\infty}a_{n}x^{n}$ 必是某个函数 $f(x)$ 的 Taylor 级数, 即
  \[
    \sum_{n=0}^{\infty}\frac{f^{(n)}(0)}{n!}x^{n}=\sum_{n=0}^{\infty}a_{n}x^{n}.
\]
  进而证明: 有无穷多个函数, 使其 Taylor 级数为任一给定的幂级数.
  \item 设函数 $f(x)=\sum_{n=0}^{\infty}a_{n}(x-x_{0})^{n}$, 且
  \[
    |n!a_{n}|<M,\quad n=1,2,\cdots,
\]
  其中 $M$ 为常数. 证明:

  (1) $f(x)$ 在任一点 $a$ 处无限次可导;

  (2) 下述展开式成立:
  \[
    f(x)=\sum_{n=0}^{\infty}\frac{f^{(n)}(a)}{n!}(x-a)^{n},\quad|x|<+\infty.
\]
\end{enumerate}
\end{exblock}
